% !TEX root = master.tex
% !TEX program = lualatex

\section{C03 — Pointers 2: complements (résumé complet)}

\subsection{Chaînes de caractères (strings) — \texttt{char[]} / \texttt{char*}}
\subsubsection{Principe fondamental}
En C, il n'existe pas de type \texttt{string}. Une chaîne est un tableau (ou pointeur sur)
\texttt{char}, terminé par le caractère nul \texttt{'\textbackslash 0'}.

\begin{lstlisting}[language=C]
char nom[6] = { 'H','e','b','u','s','\0' }; // représentation exacte
\end{lstlisting}

\subsubsection{Déclarations courantes}
\begin{itemize}
  \item Littéral : \texttt{"Bonjour"} (inclut \texttt{'\textbackslash 0'}).
  \item Tableau statique : \texttt{char s[25];} ou \texttt{char s[] = "Bonjour";}
  \item Allocation dynamique : \texttt{char* s;} (il faut \texttt{malloc/calloc}).
\end{itemize}

\begin{lstlisting}[language=C]
char nom[25];
char const welcome[] = "Bonjour";

char* dyn = NULL;
size_t n = 10;
dyn = calloc(n + 1, 1);        // +1 pour '\0'
\end{lstlisting}

\subsubsection{Affectation : piège classique}
\textbf{Erreur fréquente} :
\begin{lstlisting}[language=C]
char* s = "bonjour"; // en général PAS correct si on veut modifier s
\end{lstlisting}

\textbf{Deux cas corrects :}
\begin{enumerate}
  \item \textbf{Chaîne modifiable} : allouer puis copier (avec borne).
  \item \textbf{Chaîne non modifiable} : pointer sur littéral avec \texttt{const}.
\end{enumerate}

\begin{lstlisting}[language=C]
#include <string.h>
#include <stdlib.h>

#define TAILLE 20

// 1) modifiable
char* s = calloc(TAILLE + 1, 1);
strncpy(s, "bonjour", TAILLE);
s[TAILLE] = '\0'; // strncpy ne garantit pas toujours le '\0'

// 2) non modifiable
char const* t = "bonjour";
\end{lstlisting}

\subsubsection{Bibliothèque \texttt{<string.h>} : fonctions clés et dangers}
\begin{itemize}
  \item Copie : \texttt{strcpy} (dangereux), \texttt{strncpy} (attention au '\textbackslash 0')
  \item Concat : \texttt{strcat} (dangereux), \texttt{strncat}
  \item Compare : \texttt{strcmp}, \texttt{strncmp}
  \item Longueur : \texttt{strlen}
  \item Recherche : \texttt{strchr}, \texttt{strrchr}, \texttt{strstr}
\end{itemize}

\subsubsection{I/O de chaînes}
\begin{itemize}
  \item Écrire : \texttt{printf("\%s", s)}, \texttt{puts(s)} (ajoute \textbackslash n), \texttt{fputs(s, stdout)}
  \item Lire : \texttt{scanf("\%s", s)} (dangereux si pas borné), \texttt{fgets(s, taille, stdin)} (recommandé)
\end{itemize}

\begin{lstlisting}[language=C]
char buf[128];
fgets(buf, sizeof(buf), stdin); // lit au plus sizeof(buf)-1 et termine par '\0'
\end{lstlisting}

\subsection{Pointeurs sur fonctions}
\subsubsection{Définition et syntaxe}
Si \texttt{double f(int);} est une fonction, alors \texttt{double (*g)(int);} est un pointeur
sur une fonction de même type.

\begin{lstlisting}[language=C]
double f(int i);

double (*g)(int) = f;     // idem que &f
double z = g(3);          // idem que (*g)(3)
\end{lstlisting}

Remarque : en pratique, on n'a pas besoin d'écrire explicitement \texttt{\&} ni \texttt{*}
lors de l'affectation/appel.

\subsubsection{Passer une fonction en argument : exemple type}
\begin{lstlisting}[language=C]
typedef double (*Fonction)(double);

double integre(Fonction f, double a, double b);

double aire = integre(sin, 0.0, 3.14159);
\end{lstlisting}

\subsubsection{Exemple important : \texttt{qsort}}
Prototype :
\begin{lstlisting}[language=C]
void qsort(void* base, size_t nb_el, size_t size,
           int (*compar)(const void*, const void*));
\end{lstlisting}

Contrat de \texttt{compar} :
\begin{itemize}
  \item retourne 0 si égalité,
  \item négatif si arg1 < arg2,
  \item positif si arg1 > arg2.
\end{itemize}

\begin{lstlisting}[language=C]
#include <stdlib.h>

int compare_int(void const* arg1, void const* arg2) {
  int const* const i = arg1;
  int const* const j = arg2;
  return (*i == *j) ? 0 : ((*i < *j) ? -1 : 1);
}

int tab[100];
qsort(tab, 100, sizeof(int), compare_int);
\end{lstlisting}

\subsection{Récapitulatif : lire des déclarations avec pointeurs}
Exemples typiques (à savoir lire) :
\begin{itemize}
  \item \texttt{int* p} : pointeur sur int
  \item \texttt{int** p} : pointeur sur pointeur sur int
  \item \texttt{int (*f)()} : pointeur sur fonction retournant int
  \item \texttt{int* t[]} : tableau de pointeurs sur int
  \item \texttt{int (*p)[]} : pointeur sur un tableau d'int
\end{itemize}
Conseil : utiliser \texttt{typedef} pour rendre ça lisible.

\subsection{Cast (forçage de type)}
\subsubsection{Cast de valeurs}
\begin{lstlisting}[language=C]
double x = 5.4;
int i = (int)x; // i = 5 (tronque vers 0)
\end{lstlisting}

\subsubsection{Cast de pointeurs : danger majeur}
Caster un pointeur ne change pas le contenu mémoire, seulement l'interprétation.
\begin{lstlisting}[language=C]
double x = 5.4;
int* i = (int*)&x;

printf("%d\n", (int)x); // 5
printf("%d\n", *i);     // valeur "bizarre" (reinterpretation binaire)
\end{lstlisting}

\subsubsection{void* et casts (C vs C++)}
\begin{itemize}
  \item Vers \texttt{void*} : conversion implicite OK (C et C++).
  \item Depuis \texttt{void*} :
    \begin{itemize}
      \item en C : implicite possible,
      \item en C++ : cast explicite requis.
    \end{itemize}
\end{itemize}

\subsection{Pointeurs et tableaux}
\subsubsection{Accéder à une zone allouée comme un tableau}
\begin{lstlisting}[language=C]
double* ptr = calloc(3, sizeof(double));
// ptr pointe sur le 1er double
double a0 = ptr[0];
double a1 = ptr[1];
double a2 = ptr[2];
\end{lstlisting}

\subsubsection{Tableau statique vs tableau de pointeurs vs \texttt{int**}}
\begin{itemize}
  \item \texttt{int t[N]} \approx \texttt{int* const} vers une zone \textbf{contiguë} allouée statiquement.
  \item \texttt{int* t[N]} : tableau de N pointeurs (chaque ligne peut pointer ailleurs).
  \item \texttt{int**} : typiquement un pointeur vers une zone de pointeurs, chaque ligne allouée séparément;
        pas forcément contigu, lignes tailles variables possibles.
\end{itemize}

\subsection{Arithmétique des pointeurs}
\subsubsection{Incrémentation : dépend du type}
\begin{itemize}
  \item \texttt{p = p + 1} signifie : avancer d'un élément du type pointé.
  \item Formellement : \texttt{(char*)p + sizeof(*p)} (sans écrire ça en pratique).
  \item Pas d'arithmétique standard sur \texttt{void*}.
\end{itemize}

Parcours idiomatique :
\begin{lstlisting}[language=C]
int tab[N];
int const* const end = tab + N;
for (int* p = tab; p < end; ++p) {
  // *p
}
\end{lstlisting}

\subsubsection{Priorités : \texttt{*p++} et pièges}
\begin{lstlisting}[language=C]
char* p = s;
char lu;
while ((lu = *p++)) {
  // ici lu vaut l'ancien *p, puis p a avancé
}
\end{lstlisting}

\begin{itemize}
  \item \texttt{*p++} = \texttt{*(p++)} \to on récupere la valeur actuelle pointée puis on déplace le pointeur à l'adresse suivante.
  \item \texttt{(*p)++} incrémente la valeur pointée
  \item \texttt{*++p} On déplace d'abord le pointeur p à l'adresse suivante, puis on récupère la valeur qui s'y trouve.
\end{itemize}

\subsubsection{Soustraction de pointeurs : \texttt{ptrdiff\_t}}
\texttt{p2 - p1} donne le nombre d'éléments (du même type) entre les deux pointeurs.
Type du résultat : \texttt{ptrdiff\_t} (pas \texttt{int}).

\begin{lstlisting}[language=C]
#include <stddef.h>

ptrdiff_t d = p2 - p1;
\end{lstlisting}

\subsection{Synthèse : équivalences tableaux/pointeurs}
Pour un tableau \texttt{t} et un index \texttt{i} :
\[
  t[i] \equiv *(t+i)
\]
Et c'est \textit{symétrique} : \texttt{i[t]} existe (mais à éviter).

En paramètres de fonction :
\begin{lstlisting}[language=C]
void f(int t[N]);   // ou int t[]
// est équivalent à
void f(int* t);
\end{lstlisting}
Donc \textbf{il faut passer la taille séparément}.

\subsection{sizeof : pièges classiques}
\subsubsection{sizeof sur un tableau vs sur un pointeur}
\begin{lstlisting}[language=C]
int tab[1000];
int* t = tab;

sizeof(tab); // taille totale du tableau en octets
sizeof(t);   // taille du pointeur (4 ou 8 octets typiquement)
\end{lstlisting}

Calcul correct de N (uniquement quand on a le vrai tableau dans le scope) :
\begin{lstlisting}[language=C]
size_t N = sizeof(tab) / sizeof(tab[0]);
\end{lstlisting}

\subsubsection{Piège en fonction}
\begin{lstlisting}[language=C]
#define N 1000
void f(int t[N]) {
  // t est en réalité un int*
  // sizeof(t)/sizeof(int) -> FAUX (donne 1 ou 2 typiquement)
}
\end{lstlisting}

\subsubsection{Autre bug subtil}
\begin{lstlisting}[language=C]
int tab[1000];
const int* const end = tab + sizeof(tab); // BUG: sizeof(tab) est en octets !
for (int* p = tab; p < end; ++p) {
  utiliser(*p);
}
\end{lstlisting}
Correct :
\begin{lstlisting}[language=C]
const int* const end = tab + (sizeof(tab)/sizeof(tab[0]));
\end{lstlisting}

\subsection{Retour : flexible array member (rappel + allocation)}
Idée : structure avec un tableau « flexible » en fin pour stocker dynamiquement.

\begin{lstlisting}[language=C]
#include <stdlib.h>
#include <stdint.h>

struct vector_double {
  size_t size;
  double data[1]; // "flexible" (ancienne forme pédagogique)
};

struct vector_double* make_vec(size_t nb) {
  const size_t N_MAX = (SIZE_MAX - sizeof(struct vector_double)) / sizeof(double) + 1;
  if (nb > N_MAX) return NULL;

  struct vector_double* v =
    malloc(sizeof(struct vector_double) + (nb - 1) * sizeof(double));

  if (v == NULL) return NULL;
  v->size = nb;
  return v;
}
\end{lstlisting}

Accès : \texttt{v->data[i]}.
